\documentclass[11pt,a4paper]{article}

% -------------------------------------------------
% Packages
% -------------------------------------------------
\usepackage[utf8]{inputenc}
\usepackage[T1]{fontenc}
\usepackage[brazil]{babel}
\usepackage[a4paper,top=1.2cm,bottom=1.2cm,left=1.4cm,right=1.4cm]{geometry}
\usepackage{hyperref}
\usepackage{xcolor}
\usepackage{titlesec}
\usepackage{enumitem}
\usepackage{tabularx}
\usepackage{fontawesome5}

\hypersetup{
    colorlinks=true,
    urlcolor=black
}

\definecolor{primary}{RGB}{25,55,109}

\pagestyle{empty}
\setlength{\parindent}{0pt}

\titleformat{\section}
{\Large\bfseries\color{primary}}
{}
{0pt}
{}
[\titlerule]

\begin{document}

%%%%%%%%%%%%%%%%%%%%%%%%%%%%%%%%%%%%%%%%%%%%%%%%%%
% HEADER
%%%%%%%%%%%%%%%%%%%%%%%%%%%%%%%%%%%%%%%%%%%%%%%%%%

\begin{center}

{\Huge\bfseries Rafael Caetano da Silva}

\vspace{5pt}

{\Large Engenheiro de Software Full-Stack}

\vspace{4pt}

{\small
.NET • Java • Angular • Azure • SQL • Docker
}

\vspace{6pt}

\small

\faEnvelope\
\href{mailto:rafaelkaetanodf@gmail.com}{rafaelkaetanodf@gmail.com}
\quad|\quad
\faPhone\
(61) 99398-9820
\quad|\quad
\faMapMarker*
Brasília - DF

\vspace{3pt}

\faGithub\
\href{https://github.com/Rafael-CaetanoS}{github.com/Rafael-CaetanoS}
\quad|\quad
\faLinkedin\
\href{https://linkedin.com/in/rafaelcaetanoo}{linkedin.com/in/rafaelcaetanoo}

\end{center}

\vspace{-0.2cm}
%%%%%%%%%%%%%%%%%%%%%%%%%%%%%%%%%%%%%%%%%%%%%%%%%%
\section{Objetivo}

Atuar em projetos desafiadores como Engenheiro de Software, aplicando minha experiência em desenvolvimento de sistemas, arquitetura de software e liderança técnica para construir soluções escaláveis, contribuir com a evolução das equipes e gerar valor para o negócio por meio da tecnologia.

\section{Resumo Profissional}

Engenheiro de Software Full-Stack com experiência no desenvolvimento de aplicações corporativas utilizando .NET, Java e Angular. Atuação na construção de APIs REST, comunicação em tempo real via WebSockets com SignalR, mensageria com Azure Service Bus para processamento assíncrono e desacoplamento entre serviços, arquitetura de software, integração entre sistemas, modelagem de bancos de dados relacionais e desenvolvimento de soluções voltadas para monitoramento IoT.

Além da implementação de funcionalidades, participo da definição da arquitetura das soluções, levantamento de requisitos técnicos, planejamento das demandas, revisão de código e tomada de decisões técnicas. Atualmente exerço funções de liderança, apoiando o desenvolvimento da equipe e promovendo boas práticas de engenharia de software.

Experiência na aplicação de Clean Architecture, Domain-Driven Design (DDD), SOLID, MVC e Design Patterns para construção de soluções escaláveis, seguras e de fácil manutenção.
%%%%%%%%%%%%%%%%%%%%%%%%%%%%%%%%%%%%%%%%%%%%%%%%%%

\section{Experiência Profissional}

\textbf{Coordenador Back-end}
\hfill Ago/2025 -- Atual

Grupo Orion

\begin{itemize}[leftmargin=*,noitemsep,topsep=2pt]

\item Liderança técnica da equipe Back-end, apoiando o desenvolvimento dos profissionais e promovendo boas práticas de engenharia de software.

\item Participação no levantamento e análise de requisitos técnicos junto às áreas de negócio e produto.

\item Definição da arquitetura de novas funcionalidades utilizando Clean Architecture, Domain-Driven Design (DDD) e princípios SOLID.

\item Planejamento técnico das demandas, definição de estratégias de implementação e distribuição das atividades da equipe.

\item Desenvolvimento e evolução de APIs REST utilizando ASP.NET Core.

\item Revisão de código, validação de soluções técnicas e apoio aos desenvolvedores na resolução de problemas.

\item Tomada de decisões técnicas relacionadas à arquitetura, banco de dados, integrações, desempenho e escalabilidade das aplicações.

\item Desenvolvimento de integrações entre sistemas corporativos, serviços externos e plataformas de monitoramento IoT.

\item Estruturação e manutenção de pipelines de integração e entrega contínua utilizando Azure DevOps.

\end{itemize}

\textbf{Desenvolvedor Back-end Pleno}
\hfill Abr/2025 -- Ago/2025

Grupo Orion

\begin{itemize}[leftmargin=*,noitemsep,topsep=2pt]

\item Desenvolvimento de APIs REST utilizando ASP.NET Core seguindo os princípios de Clean Architecture e DDD.

\item Implementação de testes unitários utilizando xUnit para aumento da confiabilidade das aplicações.

\item Modelagem e manutenção de bancos PostgreSQL e MySQL.

\item Integração entre sistemas utilizando APIs REST.

\item Aplicação de SOLID, Dependency Injection e Design Patterns para desenvolvimento de soluções de fácil manutenção.

\end{itemize}


\textbf{Analista de Sistemas (Estágio)}
\hfill Dez/2023 -- Dez/2024

Infraero

\begin{itemize}[leftmargin=*,noitemsep,topsep=2pt]

\item Desenvolvimento e manutenção de aplicações corporativas utilizando ASP.NET Core, Blazor e Java.

\item Evolução de sistemas internos utilizados em operações aeroportuárias.

\item Desenvolvimento de relatórios utilizando Microsoft Report Builder.

\item Utilização de Docker para padronização dos ambientes de desenvolvimento.

\item Desenvolvimento utilizando SQL Server, Oracle, PostgreSQL e MySQL.

\item Versionamento de código utilizando Git e Azure DevOps.

\end{itemize}


%%%%%%%%%%%%%%%%%%%%%%%%%%%%%%%%%%%%%%%%%%%%%%%%%%
\section{Projetos Relevantes}

\textbf{Plataforma Cortex}

Responsável pelo desenvolvimento do backend da plataforma de monitoramento IoT, incluindo gerenciamento de alertas, notificações, configurações de alarmes e integração com dispositivos inteligentes.

\vspace{0.2cm}

\textbf{Integração Infraspeak}

Desenvolvimento de integrações entre APIs para sincronização de ativos, ordens de serviço, falhas e propriedades adicionais entre sistemas corporativos.

\vspace{0.2cm}

\textbf{SIVWEB}

Sistema responsável pelo gerenciamento de monitores aeroportuários desenvolvido utilizando ASP.NET Core e Blazor.

\vspace{0.2cm}

\textbf{SICOA}

Sistema corporativo para identificação e controle de acesso de usuários.

%%%%%%%%%%%%%%%%%%%%%%%%%%%%%%%%%%%%%%%%%%%%%%%%%%
\section{Competências Técnicas}

\textbf{Back-end}
.NET, ASP.NET Core, Java, Spring Boot

\vspace{0.15cm}

\textbf{Front-end}
Angular, JavaScript, TypeScript, HTML, CSS

\vspace{0.15cm}

\textbf{Arquitetura e Engenharia de Software}

Clean Architecture, Domain-Driven Design (DDD), MVC, Arquitetura em Camadas, SOLID, CQRS


\vspace{0.15cm}

\textbf{ORM}
Entity Framework, Dapper, Hibernate

\vspace{0.15cm}

\textbf{Banco de Dados}
SQL Server, PostgreSQL, MySQL, Oracle

\vspace{0.15cm}

\textbf{Cloud e DevOps}
Azure, Docker, Azure DevOps, Git, CI/CD

\vspace{0.15cm}

\textbf{Mensageria e Tempo Real}
Azure Service Bus, SignalR, WebSockets
\vspace{0.15cm}

\textbf{Ferramentas}
Swagger, Postman, Report Builder

%%%%%%%%%%%%%%%%%%%%%%%%%%%%%%%%%%%%%%%%%%%%%%%%%%
\section{Formação}

\textbf{Pós-graduação em Engenharia de Software}
\hfill 2025 -- 2026

Universidade Católica de Brasília

\vspace{0.2cm}

\textbf{Bacharelado em Sistemas de Informação}
\hfill 2021 -- 2024

Centro Universitário do Distrito Federal (UDF)

%%%%%%%%%%%%%%%%%%%%%%%%%%%%%%%%%%%%%%%%%%%%%%%%%%
\section{Idiomas}

Português — Nativo

Espanhol — Intermediário

Inglês — Básico

\end{document}
